\documentclass[10pt,twocolumn]{article}
\usepackage[margin=0.75in]{geometry}
\usepackage{amsmath,amssymb,booktabs,graphicx,url}
\usepackage[hidelinks]{hyperref}
\hypersetup{pdftitle={Conditioning and Directionality in Adversarial Transfer on CIFAR-10},pdfauthor={Jaiveer Bassi}}
\title{Conditioning and Directionality in Adversarial Transfer on CIFAR-10}
\author{Jaiveer Bassi\\Independent Researcher\\\texttt{jaiveerbassi@yahoo.com}}
\date{September 2026}
\begin{document}
\maketitle
\begin{abstract}
Source attack success rate (ASR) and pairwise transfer rate (PTR) are routinely reported for adversarial transfer, but neither determines transfer conditional on source success (CTR): the ASR denominator is broader than the pair-dependent population CTR requires, and both omit target behavior when the source attack fails. We derive the exact identity connecting these quantities, $\operatorname{PTR}_{st}=a_{st}\operatorname{CTR}_{st}+(1-a_{st})b_{st}$, and measure it on the full CIFAR-10 test set across three independently trained checkpoints each of ResNet-18, VGG16, and MobileNetV2. At a $2/255$ perturbation budget, the failed-source residual $b_{st}$ is at most 0.68\% in five of six model-pair directions, so CTR is almost fully determined by PTR and pairwise source success alone; VGG16-to-MobileNetV2 is the exception, with $b_{st}=2.83\%$ and a 1.19-point deviation from the zero-residual prediction. At $8/255$, transfer is strongly directional, with mean PGD PTR ranging from 18.09\% for ResNet-18-to-VGG16 to 95.23\% for VGG16-to-MobileNetV2, and this asymmetry survives a matched-learning-rate control. The result gives practitioners a way to check, from the failed-source stratum, when a reported transfer rate is arithmetic rescaling and when it reflects a genuinely separate target-side effect.
\end{abstract}

\section{Introduction}
Pairwise transfer rate and transfer conditional on source success are linked by an exact identity, $\operatorname{PTR}_{st}=a_{st}\operatorname{CTR}_{st}+(1-a_{st})b_{st}$, in which $a_{st}$ is source attack success on the jointly clean-correct population for the pair and $b_{st}$ is target error on that population when the source attack fails. The identity separates two regimes that a reported transfer number does not distinguish: when $b_{st}=0$, CTR is exactly $\operatorname{PTR}_{st}/a_{st}$ and conditioning on source success is a denominator change with no separate content, whereas when $b_{st}>0$ the successful- and failed-source strata behave differently and the conditioning carries target-side information. Measuring both terms on the full CIFAR-10 test set at $2/255$, we find $b_{st}$ at most 0.68\% in five of six directions, but 2.83\% for VGG16-to-MobileNetV2, where observed CTR falls 1.19 points below the zero-residual value---a concrete instance of the second regime inside an otherwise near-mechanical matrix.

The practical obstacle is that the identity cannot be evaluated from conventionally reported quantities. Source ASR is not $a_{st}$: its denominator counts every source-correct input, including those on which the target is already wrong, so it is target-independent where $a_{st}$ is pair-dependent. Neither ASR nor PTR records $b_{st}$ at all, because both are silent about what the target does when the source attack fails. Recovering CTR therefore requires reporting the pairwise eligible set, pairwise source success, and the failed-source residual explicitly, which is what the measurements here supply, with numerator and denominator counts for every rate.

Applying this protocol also sharpens the directional picture. At $8/255$, mean PGD PTR spans 18.09\% for ResNet-18-to-VGG16 to 95.23\% for VGG16-to-MobileNetV2 over three independently trained checkpoints per architecture, and the ordering holds in every seed. Because VGG16 required a different initial learning rate to train at all, we ran a matched-learning-rate control that retrains the other two architectures at that rate; the largest asymmetries persist, so the effect is not an artifact of the training recipe alone. The measurement protocol itself was prompted by an error in our own earlier implementation, which compared CIFAR-10 labels against 1,000-class ImageNet indices and counted pre-existing clean errors as attack successes; the audit is summarized in Section~\ref{sec:background} and detailed in the supplementary artifact.

The contributions are:
\begin{enumerate}
\item the exact decomposition $\operatorname{PTR}_{st}=a_{st}\operatorname{CTR}_{st}+(1-a_{st})b_{st}$, together with the resulting test---evaluated on the failed-source stratum---for when a reported CTR rescales PTR and when it reflects a separate target-side effect;
\item a full-test, three-seed measurement of every term in that identity, with jointly clean populations, exact counts, one direction ($b_{st}=2.83\%$) in which the residual is not negligible, and PGD PTR spanning 18.09--95.23\% at $8/255$; and
\item convergence and matched-learning-rate controls that preserve the main directional result while exposing its dependence on optimization.
\end{enumerate}
We claim none of the following: a new attack, evidence that denominator errors occur at any particular rate in the literature, or the discovery of transfer asymmetry, which is established \cite{demontis,dong}. The contribution is the decomposition and a measurement complete enough to evaluate it.

\section{Background and scope}
\label{sec:background}
Adversarial examples expose a gap between ordinary predictive accuracy and behavior under deliberately chosen perturbations. Transfer attacks use a source model to construct an input that is subsequently evaluated on another model, with the target supplying neither gradients nor feedback during construction. Transferability is therefore a property of a source, a target, an attack, and an evaluation population, and it cannot be inferred from a high target error rate alone: if the target already misclassifies an input, an unchanged input can appear to be a successful attack under an unconditional metric, and a mismatch between label spaces makes this worse.

Goodfellow et al. connect adversarial vulnerability to high-dimensional linear behavior and introduce FGSM \cite{goodfellow}. Madry et al. formulate robustness using optimization over a perturbation set \cite{madry}. Carlini and Wagner develop optimization-based attacks and demonstrate the importance of careful evaluation \cite{carlini}. Papernot et al. demonstrate practical transfer-based black-box attacks using substitute models \cite{papernot}. Liu et al. study non-targeted and targeted transfer across large models \cite{liu}, while Tramer et al. characterize shared adversarial subspaces and limits of transfer \cite{tramerspace}.

Directionality and the tension between white-box optimization and transfer are established phenomena: Demontis et al. explain asymmetric transfer through target vulnerability and surrogate complexity \cite{demontis}, while Dong et al. use momentum to improve black-box transfer \cite{dong}. We take this directional behavior as established and ask instead which population a reported transfer rate is conditioned on, and how far that choice moves the rate.

Carlini et al. systematize the methodological foundations and common pitfalls of adversarial-robustness evaluation \cite{carlini2019}. Croce and Hein introduce AutoAttack, an ensemble intended to reduce dependence on attack tuning and improve evaluation reliability \cite{croce2020}. Tramer et al. emphasize defense-specific adaptive evaluation and illustrate how apparently adequate evaluations can miss weaknesses \cite{tramer2020}. RobustBench standardizes robustness benchmarking under specified tasks and threat models and uses AutoAttack as part of its evaluation framework \cite{robustbench}.

Our study complements these evaluation efforts by examining source-to-target transfer and the populations on which transfer rates are conditioned. It does not introduce a stronger attack, certify robustness, or replace adaptive defense evaluation. The FGSM/PGD study does not constitute an AutoAttack evaluation; citing AutoAttack and RobustBench must not be interpreted as having run either benchmark. Attack-strength sensitivity and the source white-box controls are necessary for interpreting the transfer measurements.

The present protocol concerns untargeted digital attacks against deterministic image classifiers. It does not evaluate physical attacks, detection systems, adversarially trained defenses, or deployment security. CIFAR-10 results would not by themselves justify claims about autonomous vehicles or face recognition. The selected architectures are established CNN baselines, not a comprehensive representation of contemporary vision systems.

The requirements that the protocol below enforces---compatible labels, raw-pixel budgets, clean eligibility, and recorded source success---were fixed after an audit of repository commit \texttt{58f8102}, which violated each of them and whose rates and rankings are excluded from the empirical record here. That audit is a motivation for the estimands, not part of the evidence; the line-by-line version is retained in the supplementary artifact.

\section{Why error is not attack success}
Let $f_t$ be the target predictor, $y$ the true label, and $A_s(x)$ an attack constructed using source $s$. Define clean error $e_t=\Pr[f_t(x)\ne y]$, newly induced error $q=\Pr[f_t(A_s(x))\ne y\mid f_t(x)=y]$, and persistence of existing error $r=\Pr[f_t(A_s(x))\ne y\mid f_t(x)\ne y]$. By partitioning on clean correctness,
\begin{equation}
\Pr[f_t(A_s(x))\ne y]=(1-e_t)q+e_t r.
\end{equation}
For the identity attack, $q=0$ and $r=1$, so adversarial error equals clean error. An illustrative classifier with 99\% clean error therefore has 99\% adversarial error under an identity attack and zero newly induced errors. This is an algebraic example, not a measurement from CIFAR-10.

The unconditional attacked-error rate and CTR condition on different populations, so their difference is not an estimate of inflation caused by clean errors. Equation~(1) attributes all-input attacked error to newly induced and persistent clean errors; PTR and CTR then add source-correctness and source-success conditions, with empty conditioning populations remaining undefined.

\subsection{Recorded estimands}
For a fixed evaluation set of $N$ examples, let $C_s(i)$ and $C_t(i)$ indicate clean correctness. Let $S_s(i)$ and $S_t(i)$ indicate misclassification of the source-generated attacked input. Define the pairwise eligible set $E_{st}=\{i:C_s(i)\land C_t(i)\}$ and successful-source subset $F_{st}=\{i\in E_{st}:S_s(i)\}$. We report
\begin{align}
\operatorname{ASR}_s &= \frac{\sum_{i:C_s(i)} S_s(i)}{\sum_i C_s(i)},\\
\operatorname{PTR}_{st} &= \frac{\sum_{i\in E_{st}} S_t(i)}{|E_{st}|},\\
\operatorname{CTR}_{st} &= \frac{\sum_{i\in F_{st}} S_t(i)}{|F_{st}|}.
\end{align}
PTR measures target failure on jointly clean-correct inputs whether or not the source attack succeeded. CTR measures transfer conditional on source success. To expose the arithmetic effect of conditioning, define pairwise source success $a_{st}=\Pr[S_s\mid E_{st}]$ and failed-source target error $b_{st}=\Pr[S_t\mid E_{st},\neg S_s]$. Partitioning the eligible set gives the exact identity
\begin{equation}
\operatorname{PTR}_{st}=a_{st}\operatorname{CTR}_{st}+(1-a_{st})b_{st}.
\end{equation}
Consequently, $\operatorname{CTR}_{st}=\operatorname{PTR}_{st}/a_{st}$ only when $b_{st}=0$. The conventional $\operatorname{ASR}_s$ above cannot be substituted for $a_{st}$ because its denominator includes source-correct inputs on which the target was already wrong. We report $a_{st}$ and $b_{st}$ for the conditioning analysis. The diagonal is a white-box control; its CTR is tautologically one whenever its denominator is positive.

The artifact tables carry numerator and denominator counts for every reported fraction. An empty denominator is undefined, not zero. Clean accuracy and attacked accuracy on the full evaluation set are also reported. Main cross-seed results use means and seed ranges; per-run artifact tables retain Wilson 95\% intervals under a binomial sampling interpretation for fixed checkpoints. Those intervals are descriptive for the finite CIFAR-10 test set and do not measure variability across training runs. Pairwise eligibility sets differ, so raw rates across model pairs are not strictly matched-population comparisons.

\section{Experimental protocol}
\subsection{Data and model training}
The pipeline uses the standard CIFAR-10 class order \cite{cifar}. A fixed split seed of 1729 partitions the 50,000 training images into 45,000 training and 5,000 validation images. Validation data receive no random augmentation. Training uses random crops with four-pixel padding and random horizontal flips. The 10,000 test images are reserved for final evaluation; checkpoint selection uses validation accuracy only.

All networks have ten output classes and are initialized without ImageNet weights. The ResNet-18 stem uses a $3\times3$ convolution at stride one and no initial max pool. VGG16 uses its convolutional feature stack, global $1\times1$ pooling, and a single 512-to-10 linear classifier. MobileNetV2 changes its initial convolution to stride one. These adaptations, especially the compact VGG classifier, must be stated when comparing with other implementations.

Inputs to the attack interface are floating-point pixels in $[0,1]$. Channel normalization with mean $(0.4914,0.4822,0.4465)$ and standard deviation $(0.2470,0.2435,0.2616)$ occurs inside the model, so gradients include preprocessing while perturbation budgets remain in raw-pixel coordinates. Checkpoints contain architecture, class order, split indices, training seed, selected epoch, and model state. Evaluation rejects incompatible checkpoint metadata.

Training uses 200 epochs, SGD with momentum 0.9, weight decay $5\times10^{-4}$, cosine learning-rate decay, and batch size 128. ResNet-18 and MobileNetV2 use initial learning rate 0.1. That rate caused the VGG16 loss to remain near $\log 10$ and validation accuracy near chance, so we stopped that failed run and selected 0.01 using validation performance only. We then froze 0.01 for all three VGG16 seeds before test evaluation. The independently trained seeds are 0, 1, and 2; the split seed remains 1729.

\subsection{Attacks and controls}
FGSM uses one signed cross-entropy gradient step with projection into the pixel range. PGD uses random initialization within the intersection of the pixel range and the $L_\infty$ ball, 40 steps, absolute step size $2/255$, and five restarts at the primary budget $8/255$. Across iterates and restarts, the implementation prioritizes source misclassification, then larger source loss. Target predictions never influence this selection. This source-selected construction may differ from final-iterate PGD in its transfer behavior, so the selection rule is part of the attack specification.

The clean control leaves inputs unchanged. The random-noise control samples independent uniform perturbations within the same $L_\infty$ budget and clips to the pixel range. The empirical scope is limited to the $L_\infty$ threat model. We do not mix CW-$L_2$ results with these measurements or imply that equal numeric radii across norms represent equal threat models.

For every source and attack, the same generated inputs are evaluated on every target, including the source. Models remain in evaluation mode. Runtime checks reject nonfinite pixels, invalid pixel ranges, and norm violations. A saved example records its test index, true label, predictions, and actual perturbation. Visualization does not presume attack success or establish perceptual invisibility.

\subsection{Reproducibility and sensitivity analysis}
A seeded permutation specifies test indices, and the identical ordered sample list was used for every pair. All reported final evaluations used the full 10,000-image CIFAR-10 test set; the earlier 1,000-example runs were labeled as pilots and excluded from the empirical results. Each completed run saves arguments, software versions, commit identifier, checkpoint SHA-256 hashes, test indices, elapsed time, clean and attacked per-example predictions, and per-example $L_2$ and $L_\infty$ norms. Tables are generated from measured summaries rather than embedded constants.

We evaluate $L_\infty$ radii $2/255$, $4/255$, and $8/255$. The PGD step size is one quarter of the radius. At $8/255$, seed 0 uses 10, 40, and 100 steps with one restart and 100 steps with five restarts; seeds 1 and 2 repeat the 10-step, one-restart and 100-step, five-restart endpoints. At $2/255$, seed 0 also uses 100 steps and five restarts. Attack hyperparameters are not selected to maximize target transfer. We show training seeds individually and summarize seed-level estimates; we do not pool repeated test images as independent observations.

To test whether the architecture-specific VGG16 learning rate alone explains the largest asymmetries, we train seed-0 ResNet-18 and MobileNetV2 controls at 0.01 for the same 200 epochs and rerun the full-test $8/255$ protocol with the existing VGG16 seed-0 checkpoint. This is a prespecified sensitivity comparison; it is not pooled with the three-seed primary estimates.

\section{Results}
\subsection{Clean accuracy and primary attack strength}
Table~\ref{tab:clean} reports validation-selected checkpoint performance on the untouched test set. Variation across the three training seeds is small relative to the differences among architectures. At the primary $8/255$ budget, 40-step, five-restart PGD attains 100\% source ASR for every ResNet-18 and MobileNetV2 checkpoint and at least 99.98\% for every VGG16 checkpoint. The diagonal therefore confirms that the primary attack is strong against its source models.

Uniform random noise at the same $8/255$ radius produces mean off-diagonal PTR between 1.45\% and 7.89\%, below the corresponding adversarial results. Noise CTR is based on much smaller source-success subsets and is not used to rank model pairs. Because each nominal source row uses a separately seeded noise realization, these controls establish a non-adversarial baseline within each row rather than a paired comparison among noise rows.

\begin{table}[t]
\centering
\caption{Clean test accuracy (\%) by training seed. The final column is the mean and sample standard deviation over three checkpoints.}
\label{tab:clean}
\resizebox{\columnwidth}{!}{%
\begin{tabular}{lrrrr}
\toprule
Model & Seed 0 & Seed 1 & Seed 2 & Mean $\pm$ SD\\
\midrule
ResNet-18 & 95.11 & 95.01 & 95.20 & 95.11 $\pm$ 0.10\\
VGG16 & 90.59 & 90.41 & 90.88 & 90.63 $\pm$ 0.24\\
MobileNetV2 & 89.97 & 90.57 & 90.38 & 90.31 $\pm$ 0.31\\
\bottomrule
\end{tabular}
}
\end{table}

\begin{table}[t]
\centering
\caption{Mean source ASR (\%) over three seeds.}
\label{tab:asr}
\resizebox{\columnwidth}{!}{%
\begin{tabular}{lrrrr}
\toprule
Source & FGSM $8/255$ & PGD $2/255$ & PGD $4/255$ & PGD $8/255$\\
\midrule
ResNet-18 & 58.54 & 94.47 & 99.91 & 100.00\\
VGG16 & 87.31 & 71.73 & 97.93 & 99.98\\
MobileNetV2 & 92.03 & 99.85 & 100.00 & 100.00\\
\bottomrule
\end{tabular}
}
\end{table}

\subsection{Directionality at the primary budget}
Table~\ref{tab:primary} shows off-diagonal PGD PTR at $8/255$. Transfer is strongly directional. VGG16 examples transfer to ResNet-18 at 92.73\% on average, whereas the reverse direction reaches 18.09\%. MobileNetV2-to-VGG16 is 25.94\%, whereas VGG16-to-MobileNetV2 is 95.23\%. The ordering persists across all three seeds, so it is not produced by one checkpoint draw. These are associations among the trained architectures and recipes; differences in capacity, optimization, and clean accuracy preclude an architecture-only causal interpretation.

\begin{table}[t]
\centering
\caption{PGD PTR (\%) at $8/255$, reported as mean [minimum, maximum] over three training seeds.}
\label{tab:primary}
\resizebox{\columnwidth}{!}{%
\begin{tabular}{lrrr}
\toprule
Source & ResNet-18 & VGG16 & MobileNetV2\\
\midrule
ResNet-18 & --- & 18.09 [16.73,19.38] & 72.27 [71.76,72.73]\\
VGG16 & 92.73 [92.40,92.99] & --- & 95.23 [94.80,95.61]\\
MobileNetV2 & 67.60 [63.46,70.61] & 25.94 [25.73,26.29] & ---\\
\bottomrule
\end{tabular}
}
\end{table}

The matched-learning-rate control preserves the two largest VGG asymmetries (Table~\ref{tab:lrcontrol}). VGG16-to-ResNet-18 exceeds the reverse direction by 63.27 points, and VGG16-to-MobileNetV2 exceeds the reverse by 74.57 points. The control ResNet-18 has 93.82\% clean test accuracy, compared with 95.11\% for its primary seed-0 checkpoint; control MobileNetV2 has 90.90\%, compared with 89.97\%. Individual transfer entries change by up to 25.93 points across the full matrix. The control rejects a learning-rate-only account of the VGG asymmetry; it licenses no stronger causal claim because the resulting checkpoints also differ in accuracy, optimization trajectory, and loss geometry.

\begin{table}[t]
\centering
\caption{Seed-0 PGD PTR (\%) at $8/255$ before and after training all three architectures with initial learning rate 0.01.}
\label{tab:lrcontrol}
\begin{tabular}{lrr}
\toprule
Direction & Primary & Matched LR\\
\midrule
ResNet-18 $\rightarrow$ VGG16 & 19.38 & 32.71\\
VGG16 $\rightarrow$ ResNet-18 & 92.99 & 95.98\\
MobileNetV2 $\rightarrow$ VGG16 & 25.73 & 18.55\\
VGG16 $\rightarrow$ MobileNetV2 & 94.80 & 93.12\\
\bottomrule
\end{tabular}
\end{table}

\begin{table}[t]
\centering
\caption{FGSM PTR (\%) at $8/255$, reported as mean [minimum, maximum] over three seeds.}
\label{tab:fgsm}
\resizebox{\columnwidth}{!}{%
\begin{tabular}{lrrr}
\toprule
Source & ResNet-18 & VGG16 & MobileNetV2\\
\midrule
ResNet-18 & --- & 27.80 [26.21,30.01] & 58.30 [56.73,59.43]\\
VGG16 & 55.27 [54.14,56.75] & --- & 64.88 [64.38,65.71]\\
MobileNetV2 & 47.37 [47.20,47.55] & 29.34 [28.09,30.55] & ---\\
\bottomrule
\end{tabular}
}
\end{table}

\subsection{Budget and conditioning sensitivity}
Transfer falls sharply with the perturbation budget, but the directional ordering is broadly stable (Figure~\ref{fig:budget}). VGG16-to-ResNet-18 drops from 92.73\% at $8/255$ to 15.45\% at $2/255$; ResNet-18-to-MobileNetV2 drops from 72.27\% to 21.61\%. These changes rule out treating one reported transfer percentage as a property of a model pair independent of its threat model.

\begin{figure}[t]
\centering
\includegraphics[width=\columnwidth]{../artifacts/study/transfer_summary.pdf}
\caption{Mean PGD PTR across three seeds versus $L_\infty$ budget. Error bars span the minimum and maximum seed values.}
\label{fig:budget}
\end{figure}

At $2/255$, VGG16 source ASR is only 71.05--72.98\% (Table~\ref{tab:asr}). Table~\ref{tab:conditioning} applies the exact conditioning identity with absolute counts. Conditioning is nearly mechanical in five directions: mean $b_{st}$ is at most 0.68\%, and mean CTR is within 0.27 points of the value obtained by dividing mean PTR by mean $a_{st}$. Across individual seeds, the largest corresponding deviation is 0.40 points. VGG16-to-MobileNetV2 is the exception, with mean $b_{st}=2.83\%$ and observed CTR 1.19 points below the zero-residual value. This exception confirms that the successful- and failed-source strata are not interchangeable. More fundamentally, conventional source ASR and PTR cannot determine CTR even when the residual is negligible, because conventional ASR uses a broader, target-independent denominator than $a_{st}$.

\begin{table}[t]
\centering
\caption{$2/255$ PGD conditioning decomposition. Count tuples list seeds 0/1/2; rates are mean [range] percentages.}
\label{tab:conditioning}
\resizebox{\columnwidth}{!}{%
\begin{tabular}{lrr}
\toprule
Direction & Pairwise eligible $|E|$ & Source success $|F|$\\
\midrule
R$\rightarrow$V & 8905/8874/8923 & 8381/8318/8424\\
R$\rightarrow$M & 8831/8894/8873 & 8307/8339/8374\\
V$\rightarrow$R & 8905/8874/8923 & 6457/6257/6303\\
V$\rightarrow$M & 8583/8589/8620 & 6139/5976/6003\\
M$\rightarrow$R & 8831/8894/8873 & 8822/8879/8857\\
M$\rightarrow$V & 8583/8589/8620 & 8574/8574/8604\\
\bottomrule
\end{tabular}
}
\vspace{3pt}
\resizebox{\columnwidth}{!}{%
\begin{tabular}{lrrrr}
\toprule
Direction & Pair ASR & Fail-target & PTR & CTR\\
\midrule
R$\rightarrow$V & 94.09 [93.73,94.41] & 0.00 [0.00,0.00] & 5.84 & 6.21\\
R$\rightarrow$M & 94.07 [93.76,94.38] & 0.25 [0.19,0.36] & 21.61 & 22.96\\
V$\rightarrow$R & 71.22 [70.51,72.51] & 0.68 [0.42,0.96] & 15.45 & 21.42\\
V$\rightarrow$M & 70.25 [69.58,71.53] & 2.83 [2.71,2.99] & 23.29 & 31.97\\
M$\rightarrow$R & 99.85 [99.82,99.90] & 0.00 [0.00,0.00] & 10.79 & 10.81\\
M$\rightarrow$V & 99.84 [99.81,99.90] & 0.00 [0.00,0.00] & 4.61 & 4.62\\
\bottomrule
\end{tabular}
}
\end{table}

The stronger 100-step, five-restart $8/255$ endpoint changes off-diagonal PTR by at most 1.70 percentage points relative to the primary 40-step, five-restart setting across all seeds. Increasing optimization does not increase every transfer entry: for example, ResNet-18-to-VGG16 falls slightly. At $2/255$ on seed 0, increasing from 40 to 100 steps at five restarts changes all off-diagonal PTR entries by at most 0.27 points and CTR entries by at most 0.40 points. This supports convergence of the primary step count without asserting a global optimum.

FGSM reinforces the distinction between source strength and transfer (Table~\ref{tab:fgsm}). At $8/255$, its mean PTR is 27.80\% for ResNet-18-to-VGG16 and 29.34\% for MobileNetV2-to-VGG16, compared with 18.09\% and 25.94\% for PGD. PGD nevertheless has far higher source ASR. For the other four off-diagonal directions, PGD transfers more strongly than FGSM. This replicates prior observations that iterative source optimization can reduce transfer \cite{dong}; an attack ranking depends on direction and cannot be inferred from white-box source success alone.

\section{Evidence status and limitations}
The code, manuscript, and reproducibility snapshot are archived under concept DOI \href{https://doi.org/10.5281/zenodo.22737839}{10.5281/zenodo.22737839}, which resolves to the most recent version; the snapshot corresponding to this manuscript is release v1.0.2, \href{https://doi.org/10.5281/zenodo.22825090}{10.5281/zenodo.22825090}. The per-example experiment records are deposited separately under concept DOI \href{https://doi.org/10.5281/zenodo.22756878}{10.5281/zenodo.22756878}, version 1.0.1, \href{https://doi.org/10.5281/zenodo.22756879}{10.5281/zenodo.22756879}. The final record contains nine checkpoints, full-test predictions, exact eligibility counts, perturbation norms, manifests, and checkpoint and source hashes. The three seed triads provide evidence about variability under the specified training procedure, but three draws support descriptive ranges rather than precise population-level inference. Per-run Wilson intervals describe fixed-checkpoint binomial variation and do not capture training variability. The same 10,000 test images recur across seeds, so image-level counts must not be pooled as 30,000 independent observations.

The study covers three standard CNN families on CIFAR-10 and untargeted $L_\infty$ attacks. It does not establish behavior on larger datasets, transformers, adversarially trained models, targeted attacks, or physical settings. It also does not compare against every modern transfer attack. The VGG16 learning-rate adjustment was selected using validation results after the common recipe failed. A seed-0 matched-learning-rate control preserves its strongest asymmetries, but only one control seed was run and some individual entries change substantially; optimization and architecture remain entangled. Random-noise draws are independently generated for each nominal source row, so noise-control rows are reproducible but do not share one common perturbation realization. The retained manifests record a dirty working tree at execution time; they preserve prediction checksums, checkpoint hashes, the recorded commit, and hashes of the executed source files, while the clean release tag archives the final corrected code.

The original artifact is a motivating self-correction, not evidence that similar errors are common elsewhere. The present results do not imply that one architecture is intrinsically more transferable: architecture, parameterization, training dynamics, and clean accuracy vary together. PTR and CTR answer different questions, and neither should be selected after seeing which produces the more dramatic result.

\section{Conclusion}
The identity $\operatorname{PTR}_{st}=a_{st}\operatorname{CTR}_{st}+(1-a_{st})b_{st}$ fixes what a reported transfer rate does and does not determine. CTR reduces to $\operatorname{PTR}_{st}/a_{st}$ exactly when the failed-source residual vanishes, and conventional source ASR cannot stand in for $a_{st}$ because its denominator is target-independent where $a_{st}$ is pair-dependent. Measuring the residual is therefore not optional.

Nor is it always negligible. Across three training seeds at $2/255$, $b_{st}$ is at most 0.68\% in five directions, but 2.83\% for VGG16-to-MobileNetV2, where observed CTR falls 1.19 points below the zero-residual value. One direction in six is enough to show that the successful- and failed-source strata are not interchangeable in general, and that the near-mechanical cases are a finding about these pairs rather than a property of the metric.

Within the same measurements, transfer is strongly directional and budget dependent: mean PGD PTR spans 18.09--95.23\% across off-diagonal pairs at $8/255$ and 4.61--23.29\% at $2/255$. The strongest VGG directionality persists in a seed-0 matched-learning-rate control, while the movement of individual entries cautions against an architecture-only interpretation. Stronger source optimization can also reduce transfer for a particular direction, replicating established source-overfitting behavior. Compatible labels and preprocessing make these measurements meaningful; clean-error exclusions, explicit denominators, seed-level reporting, and convergence checks make their limits visible.

\begin{thebibliography}{13}
\bibitem{goodfellow} I. J. Goodfellow, J. Shlens, and C. Szegedy. Explaining and Harnessing Adversarial Examples. ICLR, 2015. \url{https://arxiv.org/abs/1412.6572}.
\bibitem{madry} A. Madry, A. Makelov, L. Schmidt, D. Tsipras, and A. Vladu. Towards Deep Learning Models Resistant to Adversarial Attacks. ICLR, 2018. \url{https://arxiv.org/abs/1706.06083}.
\bibitem{carlini} N. Carlini and D. Wagner. Towards Evaluating the Robustness of Neural Networks. IEEE Symposium on Security and Privacy, 2017. \url{https://arxiv.org/abs/1608.04644}.
\bibitem{papernot} N. Papernot, P. McDaniel, I. Goodfellow, S. Jha, Z. B. Celik, and A. Swami. Practical Black-Box Attacks against Machine Learning. ACM AsiaCCS, 2017. \url{https://arxiv.org/abs/1602.02697}.
\bibitem{liu} Y. Liu, X. Chen, C. Liu, and D. Song. Delving into Transferable Adversarial Examples and Black-box Attacks. ICLR, 2017. \url{https://arxiv.org/abs/1611.02770}.
\bibitem{tramerspace} F. Tramer, N. Papernot, I. Goodfellow, D. Boneh, and P. McDaniel. The Space of Transferable Adversarial Examples. arXiv:1704.03453, 2017. \url{https://arxiv.org/abs/1704.03453}.
\bibitem{dong} Y. Dong, F. Liao, T. Pang, H. Su, J. Zhu, X. Hu, and J. Li. Boosting Adversarial Attacks with Momentum. CVPR, 2018. \url{https://openaccess.thecvf.com/content_cvpr_2018/html/Dong_Boosting_Adversarial_Attacks_CVPR_2018_paper.html}.
\bibitem{demontis} A. Demontis, M. Melis, M. Pintor, M. Jagielski, B. Biggio, A. Oprea, C. Nita-Rotaru, and F. Roli. Why Do Adversarial Attacks Transfer? Explaining Transferability of Evasion and Poisoning Attacks. USENIX Security, 2019. \url{https://www.usenix.org/conference/usenixsecurity19/presentation/demontis}.
\newpage
\bibitem{cifar} A. Krizhevsky. Learning Multiple Layers of Features from Tiny Images. Technical report, University of Toronto, 2009. \url{https://www.cs.toronto.edu/~kriz/cifar.html}.
\bibitem{carlini2019} N. Carlini et al. On Evaluating Adversarial Robustness. arXiv:1902.06705, 2019. \url{https://arxiv.org/abs/1902.06705}.
\bibitem{croce2020} F. Croce and M. Hein. Reliable Evaluation of Adversarial Robustness with an Ensemble of Diverse Parameter-free Attacks. ICML, PMLR 119:2206--2216, 2020. \url{https://proceedings.mlr.press/v119/croce20b.html}.
\bibitem{tramer2020} F. Tramer, N. Carlini, W. Brendel, and A. Madry. On Adaptive Attacks to Adversarial Example Defenses. NeurIPS, 2020. \url{https://arxiv.org/abs/2002.08347}.
\bibitem{robustbench} F. Croce et al. RobustBench: A Standardized Adversarial Robustness Benchmark. NeurIPS Datasets and Benchmarks, 2021. \url{https://robustbench.github.io/}.
\end{thebibliography}
\end{document}
